\label{sec:kikuchi}
\subsection{What the quantum-cut sparsification paper does}
Basu, Brakensiek, Kothari and Putterman \cite{BBKP} sparsify the quantum-cut Hamiltonian
\[L_G=\textstyle\sum_ew_e(I-A_e)\quad\text{on }\mathbb C^{2^{[n]}},\]
where $A_e$ swaps the two qubits of the edge $e$, i.e.\ acts on subsets by the
transposition $(e)$. $L_G$ is the Laplacian of the Kikuchi (token) graph $K_G$ at every level $k$ at once: $S\sim T$ iff $|S|=|T|$ and
$S\oplus T\in E(G)$, the symmetric exclusion process of $k$ hard-core tokens on $G$. Their theorem: $\tilde O(n/\varepsilon^2)$ reweighted edges
preserve $L_G$ within $1\pm\varepsilon$ on every state, so one sampled graph is a spectral sparsifier of every level simultaneously.
Classical reductions are hopeless (the singlet projector has rank one, and CSPs with one satisfying assignment need $n^2$ terms);
the entangled structure is what makes it possible. The argument has four manoeuvres.
\begin{enumerate}[leftmargin=*]
\item \textbf{Up and down.} $U$ adds a token anywhere, $D$ removes one, $D=U^*$ and $[D,U]=(n-2k)I$ on level $k$. Harmonic functions are
$H_k=\ker D$ on level $k$, and $E_k=\bigoplus_{\ell=k}^{n-k}U^{(\ell)}H_k$. Each $A_e$ commutes with $U$, so every $L_G$ preserves every $E_k$, and
$L_{K_n}=k(n+1-k)$ on $E_k$.
\item \textbf{Leverage decays with the level.} On $E_k$ an edge has $L_e\preceq2$ while $L_G\succeq\lambda\frac{2m}{n^2}k(n+1-k)$, so the
matrix-Chernoff parameter grows like $k\log n/\varepsilon^2$ and beats $\log\dim E_k\lesssim k\log n$. The importance of an edge
trickles down as the level rises, and a union bound over levels closes.
\item \textbf{Up to the group, down to every level.} With $\nabla_G=\sum_ew_e(\mathrm{id}-(e))\in\mathbb C[S_n]$ and the lift $g(\sigma)=f(\{\sigma(1),\dots,\sigma(k)\})$,
$\langle f,L_ef\rangle=\langle g,\nabla_eg\rangle$: an operator inequality proved once in the group algebra holds at all levels. Alon--Kozma
give $\nabla_G\succeq\lambda\nabla_{K_n}$ for expanders, through the octopus inequality of Caputo--Liggett--Richthammer (a star at $u$
dominates the complete graph on its leaves weighted by $w_{uv}w_{uv'}/\sum w$).
\item \textbf{From expanders to everything.} Expander decomposition, Chen's lift of effective resistance to all levels,
$R_G(u,v)L_G\succeq L_{uv}$, and an edge-moving lemma across weight scales. The paper's Remark 8.5 isolates the obstruction: a bridge of
a tree keeps its full importance at every level, so the decay is an average over edges, not a per-edge fact.
\end{enumerate}

\begin{proposition}[The essence is Schur--Weyl duality]\label{prop:schurweyl}
$\mathbb C^{2^{[n]}}=(\mathbb C^2)^{\otimes n}=\bigoplus_{k\le n/2}V_{n/2-k}\otimes S^{(n-k,k)}$, where $U,D$ and $[D,U]$ generate $\mathfrak{su}(2)$ (levels are its weight
spaces), $S_n$ permutes the tensor factors, $V_j$ is the spin-$j$ representation and $S^{(n-k,k)}$ the two-row Specht module. Then
$E_k=V_{n/2-k}\otimes S^{(n-k,k)}$ and $L_G|_{E_k}=I\otimes\rho_{(n-k,k)}(\nabla_G)$. In particular every level inside $E_k$ carries the same operator, which is
why one sampling scheme serves all levels, and $L_{K_n}|_{E_k}=\binom n2-c(n-k,k)=k(n+1-k)$ with $c$ the content sum of the diagram.
\end{proposition}
\begin{proof}
$H_k=\ker D_k$ is the top constituent $S^{(n-k,k)}$ of the permutation module on $k$-subsets; $U$ commutes with $S_n$, so $U^{(\ell)}H_k\cong S^{(n-k,k)}$ for
$k\le\ell\le n-k$, and these $n-2k+1$ copies form one $\mathfrak{su}(2)$ string. $L_e=I-\rho(e)$ acts only through $S_n$. The sum of all transpositions acts on
$S^\lambda$ by the content sum; for $\lambda=(n-k,k)$ it is $\binom{n-k}2+\binom k2-k$, and $\binom n2$ minus it equals $k(n+1-k)$.
\end{proof}

\subsection{The network has the same two-duality structure, once bosonic and once hard-core}
\begin{proposition}[Howe duality of the Gaussian tokens and the radial foliation]\label{prop:howe}
On $L^2(\R^n,\gamma)$ let $E^-=\tfrac12\sum_ia_ia_i$, $E^+=\tfrac12\sum_ia_i^\dagger a_i^\dagger$, $H=N+\tfrac n2$. Then $[E^-,E^+]=H$, $[H,E^\pm]=\pm2E^\pm$, and the
triple commutes with $\Gamma(O)$, $O\in O(n)$: the pair $(O(n),\mathfrak{sl}_2(\R))$ is a Howe dual pair and
$L^2(\gamma)=\bigoplus_j\mathcal H_j\otimes M_j$, with $\mathcal H_j$ the degree-$j$ spherical harmonics (traceless token kernels, $\ker E^-$) and $M_j$
the lowest-weight $\mathfrak{sl}_2$-module of radial functions. The homogeneity operator is $x\cdot\nabla=N+2E^-$, so the Euler--Fock ladder
(Theorem~\ref{thm:ladder}) is $E^-$-descent: $E^-$ removes a pair of tokens from one vertex, where the paper's $D$ removes one hard-core
token. For a one-homogeneous $F=r\,\phi(\theta)$ the mean is the $O(n)$-invariant component, $u=\E|x|\cdot\int_{S^{n-1}}\phi\,d\sigma$: in the radial
foliation of $\R^n\setminus0$ the leafwise structure is $\mathfrak{sl}_2$, the transverse structure is $O(n)$, and the mean is the transverse
measure of the leafwise slope.
\end{proposition}
\begin{proof}
$[a_i^2,a_i^{\dagger2}]=4a_i^\dagger a_i+2$ gives $[E^-,E^+]=N+\tfrac n2$; $N$ changes by $\pm2$ under $E^\pm$; $O(n)$-invariance of $\sum a_i^2$ and $\sum a_i^{\dagger2}$
is the invariance of $\Delta$ and $|x|^2$. The decomposition is the classical theory of spherical harmonics. $x\cdot\nabla=\sum(a_i+a_i^\dagger)a_i$. The mean
formula is the polar decomposition of $\gamma$ ($|x|$ independent of $\theta$).
\end{proof}

\begin{proposition}[Gates are hard-core tokens; a layer alternates the two dualities]\label{prop:hardcore}
The gate pattern of a layer is a point $S\in2^{[n]}$, a basis vector of $(\mathbb C^2)^{\otimes n}$, and $g_a^2=g_a$ is the exclusion constraint. In the
gate-pair kernel this is exactly the contact term: $\E[g_ag_a]=\Phi_a$, where independent (bosonic) copies would give $\Phi_a^2$, so
$\diag(\Phi-\Phi^2)$ is the hard-core correction. The hidden relabelling group $S_n$ (permute the rows of $W_l$ and the columns of $W_{l+1}$) is an
exact symmetry of the network, of the He law and of every chain, and acts on gate patterns as in the paper; one realisation breaks it,
which is the odd sector of Proposition~\ref{prop:odd}. Each layer maps bosonic Gaussian tokens to hard-core gates (the selection) and
back to bosonic linear combinations (the next weights).
\end{proposition}

\subsection{What transfers: the memory has no redundancy}
\begin{proposition}[Every hub is a bridge]\label{prop:redundancy}
Let $P=\sum_cP^{(c)}$ be the readout of a source as a sum of hub atoms, and let $P_q$ keep each atom independently with probability $q$ and
weight $1/q$. Then
\[
\E\|P_q-P\|_F^2=\frac{1-q}q\sum_c\|P^{(c)}\|_F^2=\frac{1-q}{q}\,\frac{\|P\|_F^2}{\rho_{\rm eff}},\qquad\rho_{\rm eff}=\frac{\|\sum_cP^{(c)}\|_F^2}{\sum_c\|P^{(c)}\|_F^2}.
\]
On network $0$, $\rho_{\rm eff}=0.998$--$1.004$ at ages $1,3,6$ (stage-9 workflow, lens 4), so the relative error of hub sampling is $(1-q)/q$
whatever $n$. In the paper's terms every hub atom is a bridge (their Remark 8.5): its importance does not decay, because the readout
sees one level only and the atoms are orthogonal there.
\end{proposition}
\begin{proof}
The atoms are kept independently, so $\E\|P_q-P\|^2=\sum_c\Var(\xi_c/q)\|P^{(c)}\|^2$ with $\xi_c\sim\mathrm{Bernoulli}(q)$.
\end{proof}
The level decomposition that does exist on the Gaussian side (Proposition~\ref{prop:howe}) concentrates the mean in the single sector
$j=0$, the opposite of the spreading that sparsification needs. The manoeuvre that does transfer is the third: prove in the group,
descend to every level. For the network the group is the hidden relabelling $S_n$, under which the ensemble-averaged chain is symmetric;
whatever is proved there holds for all $k$-hub joint cumulants at once, and the realisation-specific (odd) part is what lies outside.

\subsection{A higher tier with random dynamics}
\begin{proposition}[Localization is a stochastic flow on the first-layer weights]\label{prop:flow}
An input $x\sim N(m,\Sigma)$ is the standard input through the effective first layer $W_1\Sigma^{1/2}$ with bias $W_1m$. Along an adapted
localization with control $C_t$, $d\Sigma_t=-\Sigma_tC_t\Sigma_t\,dt$ and $dm_t=\Sigma_tC_t^{1/2}dB_t$: the effective weights contract deterministically and the
bias is a matrix-weighted Brownian martingale. Klartag's process is the dual flow (the metric is random, the centre fixed).
\end{proposition}

\begin{theorem}[Gate martingales, wall currents and covariance consumption]\label{thm:consume}
For every layer $l$ and neuron $a$, $p_t(l,a)=\mathbb P_{\nu_t}(z_{l,a}>0)$ is a martingale with
\[
dp_t(l,a)=\big\langle J_t(l,a),\,\Sigma_tC_t^{1/2}dB_t\big\rangle,\qquad J_t(l,a)=\E_{\nu_t}\big[\delta(z_{l,a})\,\nabla z_{l,a}\big],
\]
the \emph{wall current}: the Gaussian surface integral of the wall's normal under the current ball. The conditional gate covariance
$R_t=\E_{\nu_t}[gg^{\top}]-p_tp_t^{\top}$, over all layers jointly, satisfies
\[
\E R_t=R_0-\E\int_0^t\Gamma_s\,ds,\qquad \Gamma_s=J_s^{\top}\Sigma_sC_s\Sigma_sJ_s\succeq0,\qquad R_\infty=0 ,
\]
so $\Cov(\text{gates})=\E\int_0^\infty\Gamma_s\,ds$: the contact entries $\Phi-\Phi^2$ and the exchange entries of the Kikuchi weights, within and
across layers, are integrated Grams of wall currents. The currents trickle up the layers,
$J(l,a)=\sum_bW_l(a,b)\,\E_\nu[\delta(z_{l,a})\,g_{l-1,b}\nabla z_{l-1,b}]$.
\end{theorem}
\begin{proof}
$\pi(m,\Sigma)=\E g(m+\Sigma^{1/2}Z)$ solves $\partial_\Sigma\pi=\tfrac12\nabla_m^2\pi$, so the It\^o drift $\langle\partial_\Sigma\pi,d\Sigma\rangle+\tfrac12\langle\nabla_m^2\pi,d\langle m\rangle\rangle$
vanishes, and $\nabla_m\pi=\E[\delta(z)\nabla z]$ in the sense of distributions. $\E[gg^{\top}\mid\mathcal F_t]$ is a martingale and
$d(pp^{\top})=dp\,p^{\top}+p\,dp^{\top}+d\langle p,p^{\top}\rangle$ with $d\langle p,p^{\top}\rangle=\Gamma\,dt$. Completeness of the localization and $g^2=g$ give
$R_\infty=0$. The recursion is the chain rule for $\nabla z_{l,a}$.
\end{proof}
At layer one along Eldan's path the identity is Price's (Plackett's) formula integrated along the correlation time $s=t/(1+t)$:
$\Cov(g_a,g_b)=\int_0^1\rho_{ab}\,\varphi_2(\alpha_a,\alpha_b;\rho_{ab}s)\,ds$, checked to six digits for centred and non-centred pairs and for $\rho=1$
(the contact term). This is the dynamical counterpart of Klartag's contacts. Walls whose current is non-negligible under the current
ball are in contact. The ball's gate covariance is consumed by them at the rate of their current Gram, and a wall that leaves contact
for good freezes its gate. An adaptive control $C_t$ (Lovett--Meka, Klartag) steers which walls are consumed when. By the defect-path
theorem the total error of any estimator is path-independent, but its distribution along the path is not.

\begin{theorem}[Localization resolves deep gates first]\label{thm:depthtrickle}
Two copies $X,X'$ that are conditionally independent given $\mathcal F_t$ have input correlation $t/(1+t)$. In the infinite-width limit the
expected fraction of layer-$l$ gates on which they disagree is $\theta_{l-1}(t)/\pi$, where $\theta_0=\arccos\frac t{1+t}$ and
$\cos\theta_{k+1}=J(\theta_k)/\pi$. Equivalently the neuron-averaged unresolved gate variance is $\E\Var(g_{l,a}\mid\mathcal F_t)=\theta_{l-1}(t)/(2\pi)$. It
decreases with depth, so localization freezes the deep gates before the shallow ones, strongly early ($\theta_0$ large) and weakly late
($\theta_l\approx\theta_0-l\theta_0^2/3\pi$).
\end{theorem}
\begin{proof}
$\mathbb P(g(X)\ne g(X'))=2\E[p_t(1-p_t)]$ for conditionally independent Bernoulli$(p_t)$ copies. Over the neuron ensemble the pair
$(z_{l,a}(X),z_{l,a}(X'))$ is centred Gaussian with correlation $\cos\theta_{l-1}$ (the dual-kernel recursion of Theorem~\ref{thm:folding}), and two
centred Gaussians of correlation $\cos\theta$ have different signs with probability $\theta/\pi$.
\end{proof}
\begin{center}\small
\begin{tabular}{lcccccc}\toprule
measured / predicted & layer 1 & 2 & 4 & 8 & 12 & 16\\\midrule
$t=1$ ($\rho=0.5$) & $.333/.333$ & $.287/.292$ & $.223/.236$ & $.174/.173$ & $.130/.138$ & $.136/.115$\\
$t=9$ ($\rho=0.9$) & $.140/.144$ & $.135/.136$ & $.116/.124$ & $.106/.105$ & $.082/.092$ & $.098/.081$\\
$t=99$ ($\rho=0.99$) & $.044/.045$ & $.042/.044$ & $.041/.043$ & $.040/.041$ & $.033/.039$ & $.043/.037$\\\bottomrule
\end{tabular}
\end{center}
(Fraction of disagreeing gates; He network of width $512$, $64$ pairs; the last layer drifts at finite width, as in
Proposition~\ref{prop:tokens}.)

\subsection{The two tiers form a non-commuting square}
On $L^2$ of the input and the observation noise there are two filtrations. The depth tower $\mathfrak B_1\subset\dots\subset\mathfrak B_L$ consists of the tile
algebras, generated by the gates of layers $\le l$ (the lower tier, the Kikuchi vertex sets). The localization filtration $\mathfrak F_t=\sigma(y_s,s\le t)$
is the upper tier. Each family is commutative and nested.
\begin{proposition}[Non-commuting square and the angle operator]\label{prop:square}
For $0<t<\infty$ and $l\ge1$ the pair $(\mathfrak B_l,\mathfrak F_t)$ is not a commuting square: $E_{\mathfrak F_t}$ maps a gate indicator to a smooth function strictly
between $0$ and $1$, which is not $\mathfrak B_l$-measurable, so $E_{\mathfrak B_l}E_{\mathfrak F_t}\ne E_{\mathfrak F_t}E_{\mathfrak B_l}$ and the von Neumann algebra generated by the two filtrations
is noncommutative. The angle operator $\Theta_{l,t}=E_{\mathfrak B_l}E_{\mathfrak F_t}E_{\mathfrak B_l}$ is the noise-stability Markov operator of the depth-$l$ tiling at
correlation $t/(1+t)$; its two-point marginals are the time-correlated Kikuchi weights
$K_{l,t}(a,b)=\mathbb P(g_{l,a}(X)=1,\,g_{l,b}(X')=1)$, which run from $\Phi_a\Phi_b$ at $t=0$ (independent walkers) to $K_l$ at $t=\infty$ (contacts and exchange on).
At layer one $K_{1,t}$ is the tetrachoric series with $\rho_{ab}$ replaced by $\rho_{ab}\,t/(1+t)$, the diagonal included ($\rho_{aa}=1$), so the contact
term is the resummed series at $t=\infty$. The gate process along a localization is non-Markov exactly because the square does not commute.
\end{proposition}
\begin{proof}
$\E[\1\{z>0\}\mid y_t]=\Phi(\cdot)$ of the conditional mean and variance, smooth and in $(0,1)$. For $f,g\in L^2(\mathfrak B_l)$,
$\langle f,\Theta g\rangle=\E[\E(f\mid\mathfrak F_t)\E(g\mid\mathfrak F_t)]=\E f(X)g(X')$ with $X,X'$ conditionally independent given $\mathfrak F_t$, whose correlation is
$\Var\E(X\mid y_t)=t/(1+t)$. The tetrachoric expansion of $\Phi_2$ gives the layer-one formula.
\end{proof}
Taken together: the lower-tier Kikuchi weights, conditioned on a Gaussian input law $N(m,\Sigma)$, form a field over the upper tier
that solves the heat equation, and localization paths are its Brownian characteristics. A moment chain defines a field that is
not heat-flat. Its residual is the Euler--Stein defect, and the defect-path theorem says that the error is the residual integrated
along the characteristics. This fibred picture ties both tiers together. Its leaf spaces are Hausdorff (the sphere for the radial
foliation, the space of balls for localization), so the noncommutativity of the framework lies in the non-commuting square of
Proposition~\ref{prop:square}, not in a pathological leaf space. Whether the depth tower at infinite depth, with tail equivalence of gate
sequences, yields a genuinely noncommutative tail algebra is open.
